\subsection{Contracts programming basics}

\begin{frame}[t,fragile]{From assertions to contract assertions}
\begin{itemize}
  \item Computing average grades (before C++26):
\begin{lstlisting}
double average(const std::vector<double>& grades) {
  assert(not grades.empty()); // Precondition
  double sum = 0.0;
  for (double grade : grades) {
    assert(grade >= 0.0 and grade <= 100.0); // Check
    sum += grade;
  }
  auto result = sum / grades.size();
  assert(result >= 0.0 and result <= 100.0); // Postcondition
  return result;
}
\end{lstlisting}

  \item But remember:
    \begin{itemize}
      \item \cppcode{assert} is a macro.
      \item Assertions can be disabled in production builds.
    \end{itemize}
\end{itemize}
\end{frame}

\begin{frame}[t,fragile]{Contract assertions in C++26}
\begin{itemize}
  \item C++26 introduces \textmark{contract assertions} to provide a standardized way to specify and enforce preconditions, postconditions, and invariants.
\begin{lstlisting}
double average(const std::vector<double>& grades) 
  pre(not grades.empty()) // Precondition
  post(result: result >= 0.0 and result <= 100.0) // Postcondition
{
  double sum = 0.0;
  for (double grade : grades) {       
    contract assert(grade >= 0.0 and grade <= 100.0); // Check
    sum += grade;
  }
  return sum / grades.size();
}
\end{lstlisting}
  \item Benefits of contract assertions:
    \begin{itemize}
      \item Standardized syntax and semantics for expressing contracts.
      \item Present in both debug and release builds, with violation checking depending on the selected evaluation semantics configuration.
    \end{itemize}
\end{itemize}
\end{frame}